% ============================================================================
%  Section II — Related Work  (Measurement, Elsevier)
%  Requires: booktabs, pifont, multirow. Passive voice / third person.
% ============================================================================
\providecommand{\cmark}{\ding{51}}
\providecommand{\xmark}{\ding{55}}
\providecommand{\pmark}{$\sim$}

\section{Related work}
\label{sec:related}

\subsection{Vibration-based condition monitoring of induction machines}
Vibration analysis is the most established modality for detecting mechanical and
electrical faults in induction machines, covering bearing defects, rotor
asymmetries and supply imbalance~\cite{gangsar2020signal, hamani2025advancements}.
Classical pipelines rely on industrial-grade piezoelectric accelerometers sampled
at tens to hundreds of kilohertz, which resolve the characteristic bearing defect
frequencies---ball-pass frequency outer race (BPFO), ball-pass frequency inner
race (BPFI) and ball-spin frequency (BSF)---and their harmonics through spectral
and envelope analysis. The signal models underlying that practice, and the
decomposition methods built on them, remain an active subject in their own
right~\cite{liu2026cm, li2024cm}. Vibration is not the only option: stator-current
analysis is attractive because it requires no additional transducer, and
comparative studies on induction machines and railway traction drives report that
the two modalities are complementary, with vibration generally stronger for
mechanical defects and current for electrical asymmetries~\cite{ayankoso2024cm,
sun2022cm}. Deep networks have since become the default classifier family for
both~\cite{yang2021fault, zhang2020deep, guo2024cm, chen2026cm}.

This paradigm is reflected in the public benchmarks that anchor the field,
acquired at $8$--$16$~kHz~\cite{lee2024multidomain}, $42$~kHz~\cite{sehri2024ottawa},
$51.2$~kHz~\cite{thuan2023hust} and up to $200$~kHz~\cite{huang2018bearing}.
High-rate acquisition, however, entails costly instrumentation, intrusive cabling
and limited portability, which restricts deployment to the comparatively small
population of critical machines and leaves the many low-value machines in a plant
unmonitored.

\subsection{Low-cost MEMS accelerometers as measurement instruments}
\label{subsec:mems_metrology}
Micro-electro-mechanical-systems (MEMS) accelerometers are the obvious candidate
for closing that coverage gap, and their metrological behaviour has consequently
become a subject of measurement research in its own right. The quantities that
matter are well identified: sensitivity and its traceability to a reference
standard~\cite{prato2021mems}, amplitude linearity and temperature dependence
under dynamic excitation~\cite{schiavi2023mems}, cross-axis
sensitivity~\cite{wang2025mems, najafabadi2024mems}, thermal drift and
hysteresis~\cite{xu2021mems2, kim2025mems, khankalantary2021mems}, and the
intrinsic noise of the transducer and its readout~\cite{li2020mems}. Calibration
procedures have been developed that reduce the equipment burden, from six-position
schemes~\cite{xu2021mems} to in-field methods that estimate cross-axis terms
without external references~\cite{wang2025mems} and low-cost thermal compensation
that avoids a climatic chamber~\cite{khankalantary2021mems}. Systematic
characterisations of consumer-grade devices report performance that is adequate
for many vibration tasks but markedly inferior to reference-grade transducers in
noise and bandwidth~\cite{srokosz2024mems, grouios2022accelerometers}, and the
propagation of these sensor-level limits into diagnostic features has been
examined explicitly~\cite{demilia2021mems}.

Two findings from this literature bear directly on the present study. First,
mounting is a first-order influence quantity: the position and attachment of a
MEMS accelerometer measurably changes the vibration it reports~\cite{staszewski2025mems},
which bounds how far any single-installation result can be generalised. Second,
low-cost devices are increasingly deployed as measurement instruments in their own
right, in structural and machine monitoring, once their limitations are
acknowledged and compensated~\cite{lee2026mems, sayghe2026smartphone}. Condition
monitoring practitioners have nonetheless generally regarded consumer-grade MEMS
bandwidth as insufficient for bearing diagnosis~\cite{pedotti2020lowcost}, and
purpose-built low-cost nodes therefore target kilohertz
rates~\cite{jakobsen2024lowcost, shukla2020smartsensor}.

\subsection{Smartphone-grade sensing and its measurement limitations}
Smartphones embed MEMS inertial sensors at negligible marginal cost, and their
error characteristics have been measured against laboratory references: bias,
scale-factor and noise behaviour differ appreciably between handsets and are
sensitive to acquisition settings~\cite{capuano2023smartphone}. Despite this,
smartphone acquisition has been used successfully for road-profile
estimation~\cite{zhao2019smartphone}, structural vibration monitoring with
time-synchronised handsets~\cite{periyasamy2026smartphone} and even
submillimetre acoustic vibration measurement~\cite{yang2025smartphone}. In machine
diagnosis specifically, fault detection has been demonstrated from smartphone- and
MEMS-acquired vibration using mobile applications and, more recently, deep
networks~\cite{naha2017mobile, ertargin2023deep, ertargin2025automated}. A
recurring constraint is the sampling rate: application-level acquisition is
frequently capped at or near $100$~Hz, one to three orders of magnitude below
industrial practice, so that by the Nyquist criterion the high-frequency
signatures underpinning classical bearing diagnosis are not
observable~\cite{nyquist1924certain}. What such a constrained channel can still
convey has not been established as a measurement question, and the effect of
deliberate rate reduction on diagnostic features has been examined only in
passing~\cite{he2021sampling}.

The dataset employed in this study~\cite{ertargin2026smartphone,
ertargin2026dataset} extends this line by providing long-duration,
multi-condition smartphone-MEMS recordings at a software-capped $100$~Hz.
Crucially, the accompanying baseline reports near-perfect classification accuracy
under a random split of overlapping windows, a protocol that is scrutinised in
Section~\ref{subsec:leakage}.

\subsection{Networked and edge-based condition monitoring}
The Industrial Internet of Things reframes condition monitoring as a distributed
sensing problem in which many inexpensive nodes stream, or locally distil,
diagnostic information over constrained wireless links~\cite{yousuf2024iot,
halenar2024iot}. Because raw high-rate vibration cannot be sustained over
low-power wide-area radios, a large fraction of recent work pushes feature
extraction or inference onto the node. TinyML implementations of lightweight
classifiers on microcontrollers report millisecond-scale latency and
sub-millijoule energy per inference~\cite{gao2025edge, asutkar2023tinyml,
zaidi2022tinyml}, and embedded sensor nodes now combine on-board learning with
prognostic output~\cite{chaoraingern2025tinyml}. On the acquisition side,
self-powered wireless vibration nodes~\cite{rubes2021wireless} and synchronised
long-range nodes for modal identification~\cite{qiao2024wireless} demonstrate that
the sensing, energy and communication budgets are tightly coupled. These studies
establish the feasibility of edge diagnosis, but they typically assume
industrial-bandwidth sensing and rarely quantify the trade-off between
sampling and communication cost on one side and diagnostic accuracy on the other,
which is precisely what governs an ultra-low-rate node.

\subsection{Evaluation methodology, leakage and condition shift}
\label{subsec:leakage}
A recurring threat to the credibility of learning-based fault diagnosis is data
leakage---the inadvertent sharing of information between training and test
partitions that inflates reported accuracy without improving
generalisation~\cite{kapoor2023leakage, lones2024leakage}. In vibration diagnosis
two mechanisms dominate: splitting \emph{overlapping} analysis windows drawn from
the same continuous recording, so that near-duplicate segments populate both
partitions; and recording- or component-wise correlations that let a model
recognise \emph{which recording} a window originates from rather than the fault
itself. Enforcing component-wise or recording-wise separation has been shown to
reduce headline accuracies substantially~\cite{wheat2024impact, hendriks2022towards},
and the composition of the training set has been shown to determine what a model
generalises to at all~\cite{ursu2025leakage}. The same lesson has been learned
independently in wearable sensing, where leave-one-subject-out evaluation replaced
random splitting once the latter was found to be optimistic~\cite{gholamiangonabadi2020cv}.
A related and less frequently addressed bias arises when the same
cross-validation scores are used both to select a configuration and to report its
performance; unbiased estimation requires the selection to be nested inside the
evaluation~\cite{qi2019cv}. Recent work on deep bearing diagnosis argues that
benchmark results can collapse to dataset-dependent behaviour that does not
transfer~\cite{wang2026benchmark}.

A second and complementary body of work addresses operating-condition shift, in
which training and deployment differ in speed, load or installation. Transfer and
domain-adaptation methods are the usual response~\cite{qin2021domain, zou2021domain,
huang2025domain, kaya2026domain}, and their prominence is itself evidence that
condition shift degrades performance materially. Measuring that degradation
directly, rather than attempting to remove it, is the approach taken here.
Despite these warnings, random-split evaluation remains prevalent in smartphone-
and MEMS-based machine diagnosis, including the baseline accompanying the dataset
used here~\cite{ertargin2026smartphone}.

\subsection{Reporting performance with uncertainty}
Measurement practice requires that a reported quantity carry a statement of
uncertainty, and the framework for doing so is standardised~\cite{jcgm2008gum}.
Type~A evaluation from repeated observations is well developed, including its
inverse problem~\cite{dorozhovets2020uncertainty}, its behaviour when random noise
and systematic effects act simultaneously~\cite{dorozhovets2025uncertainty}, and
its estimation from a single measurement through
propagation~\cite{ou2024uncertainty}. Extending this discipline to quantities
produced by learned models is comparatively recent: methodologies have been
proposed for GUM-conformant uncertainty evaluation of machine-learning-based
measurement models~\cite{angrisani2025uncertainty}, and for making decisions in
the presence of measurement uncertainty~\cite{morello2020uncertainty}. Diagnostic
accuracy, by contrast, is still routinely reported in the fault-diagnosis
literature as a bare point estimate, without an interval, without a statement of
the unit over which it was replicated, and often without the information needed to
reproduce the partitioning that produced it. That gap between measurement practice
and diagnostic reporting practice motivates the protocol developed in
Section~\ref{sec:method}.

\subsection{Positioning}
Table~\ref{tab:positioning} situates this work against representative prior art.
Existing studies advance one axis at a time. High-rate benchmarks and deep models
maximise accuracy but ignore cost and evaluation rigour; methodology papers expose
leakage but on industrial-grade data; MEMS metrology papers characterise the
sensor thoroughly but stop short of the diagnostic task; smartphone-based studies
embrace low-cost sensing but retain optimistic random-split evaluation; and
IoT/TinyML works demonstrate edge deployment yet assume industrial bandwidth. To
the best of the authors' knowledge, no prior study jointly evaluates ultra-low-rate
smartphone-MEMS machine diagnosis under a leakage-free, cross-condition protocol,
decomposes the resulting optimism into its mechanisms, and maps the
accuracy--resource trade-off into communication and edge-model terms. This paper
addresses that combination, and reports the result with an interval whose sampling
unit is stated.

% ---------------------------------------------------------------------------
\begin{table*}[t]
\centering
\caption{Positioning of this work against representative prior studies.
\cmark~= addressed, \pmark~= partially addressed, \xmark~= not addressed.}
\label{tab:positioning}
\renewcommand{\arraystretch}{1.25}
\begin{tabular}{@{}p{3.9cm}p{2.6cm}ccccc@{}}
\toprule
\textbf{Representative work(s)} & \textbf{Sensor / rate} &
\textbf{Low} & \textbf{Leakage-free} & \textbf{Acc.--resource} &
\textbf{IoT/edge} & \textbf{Sensor} \\
 & & \textbf{cost} & \textbf{evaluation} & \textbf{trade-off} &
\textbf{analysis} & \textbf{charact.} \\
\midrule
High-rate vibration FD \& deep models~\cite{lee2024multidomain,sehri2024ottawa,thuan2023hust,zhang2020deep,guo2024cm}
 & Industrial accel., $8$--$200$\,kHz & \xmark & \xmark & \xmark & \xmark & \xmark \\
Evaluation / leakage studies~\cite{kapoor2023leakage,wheat2024impact,hendriks2022towards,lones2024leakage,qi2019cv}
 & Industrial accel., k\,Hz & \xmark & \cmark & \xmark & \xmark & \xmark \\
MEMS metrology~\cite{prato2021mems,schiavi2023mems,xu2021mems,wang2025mems,staszewski2025mems}
 & MEMS, bench setup & \cmark & n/a & \xmark & \xmark & \cmark \\
Smartphone / mobile machine FD~\cite{naha2017mobile,ertargin2023deep,ertargin2025automated}
 & Smartphone MEMS, $\leq$\,k\,Hz & \cmark & \xmark & \xmark & \pmark & \xmark \\
Source-dataset baseline~\cite{ertargin2026smartphone}
 & Smartphone MEMS, $100$\,Hz & \cmark & \xmark & \xmark & \xmark & \pmark \\
Low-cost MEMS sensor nodes~\cite{pedotti2020lowcost,jakobsen2024lowcost,rubes2021wireless}
 & MEMS node, k\,Hz & \cmark & \xmark & \pmark & \cmark & \pmark \\
TinyML edge fault diagnosis~\cite{gao2025edge,asutkar2023tinyml,chaoraingern2025tinyml}
 & (Lab) accel., k\,Hz & \cmark & \xmark & \pmark & \cmark & \xmark \\
\midrule
\textbf{This work} & \textbf{Smartphone-grade MEMS, $100$\,Hz}
 & \cmark & \cmark & \cmark & \pmark & \pmark \\
\bottomrule
\end{tabular}
\end{table*}
